\subsection{Exact implicit output for polynomial boundary optima}\label{app:boundary}

This subsection makes precise what exact output means for the explicit
polynomials of Section~\ref{sec:polynomial} when the minimizer can be
irrational and lie on the boundary. Throughout, $F$ is an explicit rational
polynomial of fixed degree $d\ge2$ with factor scopes in the bags, on a mixed
box $X$, with a verified rational bound
$L\ge\max\{0,\max_i\sup_{\bar X}\partial_{ii}F\}$,
for example the largest bound of Lemma~\ref{lem:intcurv}. Fix rationals
$C_2,C_3\ge1$ with
\[
\sum_{j\in I_C}\sup_{\bar X}|\partial_{ij}F|\le C_2,\qquad
\sum_{j,k\in I_C}\sup_{\bar X}|\partial_{ijk}F|\le C_3
\qquad(i\in I_C).
\]
Monomialwise interval bounds give such numbers with bit length $\poly(I)$
at fixed $d$.

In this subsection every coordinate uses the single curvature bound $L$, that
is, $L_i:=L$ for all $i$; these are valid upper coordinate curvatures. If
$L>0$, then $P$ contains every coordinate, so every coordinate is gridded and
filtered, and the localization bound (G4) of Lemma~\ref{lem:commonmesh}
applies to every coordinate. If $L=0$, the first stage uses endpoint grids,
is exact and has zero gap. We use TRIAL (Algorithm~\ref{alg:trial}) with a
common mesh $h_j=s2^{-j}$ as in Lemma~\ref{lem:commonmesh}. Under point
growth~\eqref{eq:growth} with $8\kappa\theta^2\le1$, (G4) places every
retained interval of stage $j$ within
$4.2\sqrt{n\kappa}\,h_j\le5\sqrt{n\kappa}\,h_j$ of $(y_j)_i$, plus one for
integer coordinates; we use the radius $5\sqrt{n\kappa}\,h_j$ below.

\paragraph{Three sound tests.}
The trials are run with one extra filter.

\begin{lemma}[Integer-label filter]\label{lem:labelfilter}
(i) Suppose at some stage every grid interval of an integer coordinate $i$
has length one. Then deleting the nodes $v$ with $m_i(v)>U$ and taking the
integer hull of the remaining nodes removes no point with objective at most
$U$, and keeps the incumbent and the corrected minimizer. The stages remain a
valid path certificate (Definition~\ref{def:cert}): the second alternative of
condition~(C1) holds for every grid interval of coordinate $i$ that the step
removes.
(ii) Assume a unique minimizer $x^*$, point growth \eqref{eq:growth}, and
$8\kappa\theta^2\le1$. At every stage with $h_j\le1/(10\sqrt{n\kappa})$,
part (i) applies to every integer coordinate and fixes it to $x^*_i$.
\end{lemma}

\begin{proof}
(i) Since every grid interval of coordinate $i$ has length one, the integers
of the current interval are exactly the nodes. For feasible $x$ in the current
box with $x_i=v$, Proposition~\ref{prop:cellwise}(c) gives $F(x)\ge m_i(v)$.
Hence a deleted node
carries no point with value at most $U$, in particular no minimizer and not
the incumbent. The corrected minimizer $y$ survives because
$m_i(y_i)\le Q(y)=\beta\le\OPT\le U$. A removed grid interval of
coordinate $i$ has length one, hence effective width zero, and each of its
endpoints outside the new interval is a deleted node $v$ with
$m_i(v)>U\ge\beta$; this is the second alternative of (C1).

(ii) If $j=0$, then $s=h_0<1$ and there is no integer coordinate. Otherwise,
by (G4) at stage $j-1$, the current integer interval lies within
$1+5\sqrt{n\kappa}\,h_{j-1}=1+10\sqrt{n\kappa}\,h_j\le2$ of the center
$c_i=(y_{j-1})_i$. All steps $\max\{1,\lfloor h_j+\theta t\rfloor\}$ with
$t\le2$ equal one, because $h_j+\theta t\le\frac1{10}+\frac12<1$, so (i)
applies. A retained node $v$ has $m_i(v)\le U$, so the grid point that attains
$m_i(v)$ has $Q\le U$, and (G2) gives
$\abs{v-x^*_i}^2\le\frac{17}{15}\kappa nh_j^2<1$; so $v=x^*_i$. The extra
filter keeps $x^*$ and the centers and only shrinks boxes; the proof of
Lemma~\ref{lem:commonmesh} uses nothing else about the boxes, so its bounds
remain valid for the modified trial.
\end{proof}

\begin{lemma}[Monotone face reduction]\label{lem:monotone}
Let $X'$ be a product box with fixed integer coordinates, $i$ a free
continuous coordinate whose original bound $\ell_i$ is an endpoint of $X'_i$,
and suppose $\partial_iF\ge0$ on $X'$. Then
$\min_{X'}F=\min_{X'\cap\{x_i=\ell_i\}}F$, and if $\partial_iF>0$ on $X'$,
every minimizer over $X'$ has $x_i=\ell_i$. The same holds for $u_i$ with
$\partial_iF\le0$. With $c$ the midpoint and $r$ the largest half-width of the
free continuous coordinates of $X'$, $\partial_iF(c)-C_2r>0$ certifies
$\partial_iF>0$ on $X'$, and $\partial_iF(c)+C_2r<0$ certifies
$\partial_iF<0$.
\end{lemma}

\begin{proof}
Replacing $x_i$ by $\ell_i$ keeps a point of $X'$ in $X'$ and, by the mean
value theorem along the segment, does not increase $F$, and decreases it under
a strict sign unless $x_i=\ell_i$. For the test,
$|\partial_iF(x)-\partial_iF(c)|\le\sum_j\sup|\partial_{ij}F|\,|x_j-c_j|\le C_2r$.
\end{proof}

Any other sound sign enclosure may be used as well, for example Bernstein
coefficient bounds \cite{Garloff1986} of each factor derivative on its own
scope, which have at most $d^p$ coefficients. Such monotonicity tests are
standard in interval global optimization
\cite{Hansen1980,HansenWalster2004,ArayaTrombettoniNeveu2010}; the
termination proof below uses only the midpoint test.

\begin{lemma}[Uniform convexity patch]\label{lem:patch}
Let $X'$ be a box with fixed integer coordinates and some fixed continuous
coordinates, $J_f$ its free coordinates, $c$ its midpoint and $r$ the largest
free half-width. If $H_{J_fJ_f}(c)-(C_3r+\sigma)I\succ0$ for some
$\sigma>0$, then $F$ restricted to $X'$ is strongly convex in $x_{J_f}$ and has a
unique minimizer over $X'$.
\end{lemma}

\begin{proof}
Each entry satisfies $|H_{jk}(x)-H_{jk}(c)|\le r\sum_l\sup|\partial_{jkl}F|$,
so every row sum of $|H_{J_fJ_f}(x)-H_{J_fJ_f}(c)|$ is at most $C_3r$. The spectral
norm of a symmetric matrix is at most its largest absolute row sum. Hence
$H_{J_fJ_f}(x)\succ\sigma I$ on $X'$.
\end{proof}

\paragraph{The procedure.}
Trial $\mu\ge2$ runs TRIAL with a common mesh, grading $\theta=2^{-\mu}$,
the cap $8\theta^{-1}\lceil\log_2(n+2)\rceil$ of Lemma~\ref{lem:commonmesh}(G5)
and the filter of Lemma~\ref{lem:labelfilter}, for stages $j=0,1,2,\ldots$
without a stage limit, aborting if the cap would be exceeded; step~(iii) of
TRIAL is replaced by the following test. After stage $j$ the trial succeeds
if $U=\beta_j$ (a zero certified gap; the incumbent is then an exact rational
minimizer). Otherwise it filters, and if every integer coordinate of the
retained box $X^{(j+1)}$ is a singleton, it copies $X^{(j+1)}$, applies sound
reductions of Lemma~\ref{lem:monotone} to the copy until none applies, and
succeeds if no free coordinate remains or if the test of
Lemma~\ref{lem:patch} holds with $\sigma_\mu=L/4^{\mu+1}$. The copy never
influences later filtering. The procedure runs phases $k=2,3,\ldots$; in
phase $k$ it restarts trial $\mu$ for each $2\le\mu\le k$ and stops it after
$2^{k-\mu}$ bit operations, stopping at the first success.

\begin{theorem}[Exact implicit boundary output with a complementarity margin]\label{thm:boundary}
(a) On success, the output is either the incumbent after a zero certified
gap, which is then a rational minimizer, or a fixed integer assignment, a set
of continuous coordinates fixed at original bounds, and a rational box for the
remaining coordinates on which $F$ is strongly convex; in the second case the
unique minimizer $\hat x$ of $F$ over this face box is a global minimizer of
the original problem, and if no coordinate is free, $\hat x$ is an explicit
rational minimizer. This holds with no assumption on the instance.

(b) Suppose $F$ has a unique minimizer $x^*$ with point growth
\eqref{eq:growth}, and every coordinate in $A=J_\partial\cap I_C$, the
continuous coordinates at which $x^*$ is at a bound, is strictly complementary:
its inward derivative, $\partial_iF(x^*)$ at a lower bound or
$-\partial_iF(x^*)$ at an upper bound, is positive. Let $\lambda_A$ be the
smallest such inward derivative and
$B_\lambda=\lceil\log_2\max\{2,C_2/\lambda_A\}\rceil$ ($B_\lambda=1$ if
$A=\emptyset$). Then the procedure succeeds after
$f_d'(p,\kappa)\poly(I+B_\lambda)$ bit operations, for a computable function
$f_d'$. Neither $x^*$, $g$, nor $\lambda_A$ is supplied.
\end{theorem}

\begin{proof}
(a) A zero certified gap makes the incumbent optimal
(Theorem~\ref{thm:certificate}). Otherwise the retained box contains every
minimizer (Proposition~\ref{prop:filter} and Lemma~\ref{lem:labelfilter}(i)),
and its integer coordinates are singletons, so its minimum is $\OPT$. Each
reduction preserves the minimum over the current box
(Lemma~\ref{lem:monotone}). On the final face box $F$ has minimum $\OPT$ and,
by Lemma~\ref{lem:patch} or because it is a single point, a unique minimizer,
which is therefore globally optimal.

(b) If $L=0$, trial $2$ succeeds at stage $0$ with zero gap, and the phase
accounting at the end of this proof applies with $\mu_*=2$ and $j_*=0$;
assume $L>0$.
Let $\mu_*$ be the least $\mu\ge2$ with $4^\mu\ge8\kappa$, so that
$4^{\mu_*}\le32\kappa$ and trial $\mu_*$ satisfies $8\kappa\theta^2\le1$; by
Lemma~\ref{lem:commonmesh}(G5) it never reaches its cap. Let $j_*$ be the least
stage with
\[
h_j\le\frac1{10\sqrt{n\kappa}},\qquad
5\sqrt{n\kappa}\,h_j\le\frac{\lambda_A}{4C_2},\qquad
5\sqrt{n\kappa}\,h_j\le\frac L{8\cdot4^{\mu_*}C_3},
\]
omitting the first condition without integer coordinates and the second if
$A=\emptyset$. Then $j_*=\poly(I)+O(\log\kappa+B_\lambda)$, since
$\log s$, $\log C_2$, $\log C_3$ and $\log(1/L)$ are polynomial in $I$. We
show that trial $\mu_*$ succeeds by stage $j_*$.

At stage $j_*$ the integer coordinates are fixed to $x^*$
(Lemma~\ref{lem:labelfilter}(ii)) and every continuous retained interval has
half-width $r\le5\sqrt{n\kappa}\,h_{j_*}$ (Lemma~\ref{lem:commonmesh}(G4)). A
sound reduction never removes $x^*$: if a reduction that fixes coordinate $i$
at a bound $\xi_i\in\{\ell_i,u_i\}$ is justified on a box $X'\ni x^*$, then
setting $x^*_i$ to $\xi_i$ gives a point of $X'$ with value at most $\OPT$,
which by uniqueness is $x^*$; so $x^*_i=\xi_i$ and $i\in A$. Hence the
current face box always contains $x^*$, agrees with it on fixed coordinates,
and has
$\norm{c-x^*}_\infty\le r$. For $i\in A$ at $\ell_i$ that is not yet
fixed, $\ell_i$ is an endpoint of the current interval and
$\partial_iF(c)-C_2r\ge\lambda_A-2C_2r\ge\lambda_A/2>0$, so the midpoint test
fixes it; the upper case is symmetric. When no reduction applies, the free
set $J_f$ consists exactly of the continuous coordinates that are interior at
$x^*$. If $J_f=\emptyset$ the trial succeeds. Otherwise $x^*\pm tv\in X$ for
$v\in\R^{J_f}$ and small $t$, so \eqref{eq:growth} and Taylor's formula give
$\nabla_{J_f}F(x^*)=0$ and $H_{J_fJ_f}(x^*)\succeq2g\,I$. By the row-sum bound,
$H_{J_fJ_f}(c)\succeq(2g-C_3r)I$, and
\[
H_{J_fJ_f}(c)-(C_3r+\sigma_{\mu_*})I\succeq(2g-2C_3r-\sigma_{\mu_*})I
\succ0,
\]
because $2C_3r\le L/(4\cdot4^{\mu_*})\le g/32$ and
$\sigma_{\mu_*}=L/(4\cdot4^{\mu_*})\le g/32$. Exact symmetric elimination
detects this.

Through stage $j_*$, trial $\mu_*$ performs the table work of
Lemma~\ref{lem:dp} with the cap $K=8\cdot2^{\mu_*}\lceil\log_2(n+2)\rceil$ on
numbers of bit length $\poly(I+j_*+\mu_*K)$ at fixed $d$
(Section~\ref{sec:polynomial}), plus at most $n^2$ sign tests and one
elimination per stage. Call this number of bit operations $M_*$; it is at
most a computable function of $(p,\kappa)$ times $\poly(I+B_\lambda)$. Phase
$k$ allots $\sum_{\mu=2}^k2^{k-\mu}<2^{k-1}$ operations, and trial $\mu_*$
completes in phase $k_*=\mu_*+\lceil\log_2M_*\rceil$. The total is at most
$2^{k_*}\le2\cdot2^{\mu_*}M_*=O(\sqrt\kappa\,M_*)$, plus clock overhead
polylogarithmic in this bound, which proves~(b).
\end{proof}

The output is an exact implicit description, not a rational coordinate claim:
for $(x^2-\frac12)^2+y+xy^2$ on $[0,1]^2$ the face is $y=0$ and the minimizer
is $(1/\sqrt2,0)$. Weak reductions preserve the minimum but not necessarily
every minimizer, so without uniqueness the certificate proves global
optimality of $\hat x$, not uniqueness in $X$. Trials must be dovetailed: a
trial with $4^\mu<8\kappa$ may run forever without exceeding its cap, and
increasing $\mu$ once per unsuccessful stage would make the grading depend on
$\lambda_A$.

The margin term and strict complementarity are both needed for certificates
of this form, as the next two propositions show.

\begin{proposition}[The margin cost is intrinsic to rectangular sign certificates]\label{prop:margin}
For $n\ge1$ let $X=[0,\frac12]^{n+2}$ with variables $x_0,\ldots,x_n,y$, and
\[
\rho_0=x_0-\tfrac14,\quad \rho_i=x_i-\tfrac14x_{i-1}^2,\quad
\chi=x_n-\tfrac18x_{n-1}^2,\quad
G_n=\sum_{i=0}^n\rho_i^2+y^2+3y\chi.
\]
Then $G_n$ has degree four, a path decomposition with bags of size at most
three, $L=\frac{35}{16}$, the unique minimizer $(a,0)$ with
$a_i=2^{-(4\cdot2^i-2)}$, $\OPT=0$, point growth with $g=\frac9{64}$
($\kappa=\frac{140}9$), and strict complementarity at $y$ with
$\lambda_A=\frac32a_n$. Every success of the procedure on $G_n$ produces a
rational number with reduced denominator at least $2^{4\cdot2^n-2}$.
Consequently a successful trial $\mu$ reaches a stage $j$ with
$j+\mu K\ge4\cdot2^n-2-\log_2\Gamma_X$, where $K$ is its cap and
$\Gamma_X=2^{n+2}$ is the product of the denominators of the box endpoints.
\end{proposition}

\begin{proof}
The recursion $a_0=\frac14$, $a_i=\frac14a_{i-1}^2$ gives the stated $a_i$,
all in $(0,\frac12)$, and $\rho(a)=0$. With $z=x-a$, $\rho_0=z_0$ and
$\rho_i=z_i-b_iz_{i-1}$ with $b_i=(x_{i-1}+a_{i-1})/4\in[0,\frac14]$, so
$\norm\rho\ge\frac34\norm z$ and $\Phi_n=\sum_i\rho_i^2\ge\frac9{16}\norm{x-a}^2$.
Since $\chi=\frac12(x_n+\rho_n)$,
\[
G_n=\tfrac14(\Phi_n+y^2)+\tfrac34(\Phi_n-\rho_n^2)+\tfrac34(\rho_n+y)^2+\tfrac32yx_n ,
\]
and every term is nonnegative on $X$. Thus
$G_n\ge\frac9{64}(\norm{x-a}^2+y^2)$, with equality to zero only at $(a,0)$.
The diagonal second derivatives are $2+\frac34x_i^2-x_{i+1}\le\frac{35}{16}$
for $i<n-1$, the same minus $\frac34y$ for $i=n-1$, and $2$ for $x_n$ and
$y$. The bags are $\{x_{i-1},x_i\}$ and $\{x_{n-1},x_n,y\}$. Finally
$\partial_yG_n(a,0)=3\chi(a)=3(a_n-\frac12a_n)$.

Consider a success. A zero-gap success outputs the unique minimizer, whose
coordinate $a_n=2^{-(4\cdot2^n-2)}$ has the stated denominator. Otherwise
the integer part is empty, and by the argument in the proof of
Theorem~\ref{thm:boundary}(b) no sound reduction fixes an $x$-coordinate,
because all are interior at the unique minimizer. If $y$ is not fixed, the
patch test needs $H_{J_fJ_f}\succ0$ at $(a,0)$, which lies in the face box; but
the principal minor on $\{x_n,y\}$ is
$\det\bigl(\begin{smallmatrix}2&3\\3&2\end{smallmatrix}\bigr)=-5$. So $y$ is
fixed, and not at its upper bound, since that would need
$\partial_yG_n\le0$ at $(a,0)$. So it is fixed by a sound certificate of
$\partial_yG_n\ge0$ on a retained box
$X'=\prod_i[\alpha_i,\alpha'_i]\times[0,\alpha'_y]$ that contains $(a,0)$. At
its point with $y=0$, $x_n=\alpha_n$ and $x_{n-1}=\alpha'_{n-1}$,
$\partial_yG_n=3(\alpha_n-(\alpha'_{n-1})^2/8)$, so
$\alpha_n\ge(\alpha'_{n-1})^2/8\ge a_{n-1}^2/8=a_n/2>0$; also $\alpha_n\le a_n$.
A rational in $(0,2^{-(4\cdot2^n-2)}]$ has reduced denominator at least
$2^{4\cdot2^n-2}$. As in Appendix~\ref{app:growth}, with $h_j=s2^{-j}$, box
endpoints and grid nodes of stage $j$ in trial $\mu$ have denominators
dividing $\Gamma_X2^{j+\mu K}$, which proves the last claim.
\end{proof}

Here $I=O(n\log n)$ while every successful certificate of the above form
needs exponentially many bits in $n$, at fixed $d$, $p$ and $\kappa$. Thus
the $B_\lambda$ term of Theorem~\ref{thm:boundary} cannot be removed without
a different kind of certificate, for instance one that keeps the correlation
between $x_{n-1}$ and $x_n$, such as the displayed identity.

\begin{proposition}[Weak complementarity defeats sign and patch tests]\label{prop:weakcompl}
On $X=[\frac12,1]\times[0,1]^2$ with variables $(x,w_1,w_2)$ let
$u=x^2-\frac12$ and
\[
\Psi=u^2+w_1^2+w_2^2+3w_1w_2+u(w_1-w_2)
=\tfrac12(u^2+w_1^2+w_2^2)+\tfrac12(u+w_1-w_2)^2+4w_1w_2 .
\]
Then $\Psi$ has degree four and one bag of size three; its unique minimizer
is $x^*=(1/\sqrt2,0,0)$ with $\OPT=0$; point growth holds with $g=\frac12$ and
$L=12$ ($\kappa=24$); and $\nabla_w\Psi(x^*)=0$, so both active coordinates
are weakly complementary. For every rational box $X'\subseteq X$ containing
$x^*$, none of $\partial_x\Psi$, $\partial_{w_1}\Psi$, $\partial_{w_2}\Psi$ has
constant sign on $X'$, and $\nabla^2\Psi(x^*)$ is not positive semidefinite.
Hence no trial of the procedure ever succeeds on $\Psi$.
\end{proposition}

\begin{proof}
The identity is direct expansion. On $X$ it gives
$\Psi\ge\frac12(u^2+\norm w^2)$, and
$u^2=(x-x^*_1)^2(x+x^*_1)^2\ge(x-x^*_1)^2$ because $x+x^*_1>1$. So
$\Psi\ge\frac12\norm{(x,w)-x^*}^2$, with zero only at $x^*$. Also
$\partial_x^2\Psi=12x^2-2+2(w_1-w_2)\le12$ and $\partial_{w_i}^2\Psi=2$. The
$x$-interval of $X'$ contains the irrational $x^*_1$ in its interior. On the
segment $w=0$ inside $X'$, $\partial_x\Psi=4xu$, $\partial_{w_1}\Psi=u$ and
$\partial_{w_2}\Psi=-u$, and $u$ changes sign at $x^*_1$. So no coordinate can
be fixed by a sound test, and the patch test on all three coordinates would
need $\nabla^2\Psi(x^*)\succ0$; but $v=(0,1,-1)$ gives
$v^{\T}\nabla^2\Psi(x^*)v=2+2-6=-2$. A zero certified gap is impossible
because $x^*$ is irrational.
\end{proof}

The identity in Proposition~\ref{prop:weakcompl} is itself a short global
certificate that $\OPT=0$ and that $x^*$ is the unique minimizer.
Certificates that keep such correlations, rather than certify signs over
independent coordinate rectangles, are the natural route to exact output
under point growth alone; no general construction of them with parameterized
size is known.
